%------------------------------------------------------------------------------%
\chapter{Apresentação}

%------------------------------------------------------------------------------%
\section{Sobre a Disciplina}

A disciplina de \emph{Matemática 01} é composta de conteúdos comuns em disciplinas de \emph{Cálculo 01} e 
de \emph{Geometria Analítica}, com uma enfase desejável em aplicações em Administração e Economia.
Como não existe um livro que corresponda a ementa dessa disciplina, comecei a escrever essas notas, 
com o intuito de facilitar o estudo por parte dos alunos. 

\begin{observacao}
    Essas notas estão em construção e provavelmente estão incompletas e contém erros. \\
    Elas não são um substituto para as aulas.
\end{observacao}

Essas notas começam com uma revisão de conteúdos considerados do Ensino Médio (Capítulo~\ref{cap:revisao});
apresentam uma introdução a matrizes e sistemas lineares (Capítulo~\ref{cap:matrizes});
em seguida definem o plano cartesiano as equações de formas geométricas como retas e circunferências (Capítulo~\ref{cap:equacoes}).
Nesse ponto há uma mudança de enfoque com a disciplina, se voltando para o conteúdo do \emph{Cálculo 01}.
Começamos com uma revisão e aprofundamento do conceito de função (Capítulo~\ref{cap:funcoes});
depois é apresentado um conceito novo, o Limite (Capítulo~\ref{cap:limites});
tendo dominado o conceito de limite é possível definir a variação instantânea de uma grandeza, ou sua Derivada (Capítulo~\ref{cap:derivadas}) e
a disciplina termina com a apresentação de algumas Aplicações da Derivada (Capítulo~\ref{cap:aplicacoes}).
No final dessas notas estão listados os livros que foram usados como base para sua construção. 
Se possível, é interessante que o aluno utilize também essa bibliografia assim como grande quantidade 
de material de apoio encontrado na \emph{internet}, como vídeo aulas.

%------------------------------------------------------------------------------%
\section{Sobre as Avaliações}

Os alunos precisam prestar atenção a um fato muito importante sobre as avaliações da disciplina de \emph{Matemática~01}:
elas são muito diferentes das provas do ENEM.
Muitos alunos foram treinados, durante seus estudos do Ensino Médio, para responder questões
no estilo específico do ENEM. Um problema com essa abordagem é que essas questões são 
sempre de múltipla escolha, onde basta chegar a resposta correta. Ninguém 
vai avaliar ou mesmo ler os passos utilizados pelo aluno para chegar a resposta.
As avaliações dessa disciplina serão sempre discursivas, isso é, \textbf{o aluno deve escrever
as respostas de maneira completa, correta, legível e compreensível} para o professor que 
for corrigir as avaliações.

Para fazer isso o aluno deve ter atenção a vários detalhes importantes
\begin{itemize}
    \item Escrever as justificativas logicamente encadeadas
    \item Usar a notação matemática correta
    \item Estar atento ao significado de cada símbolo matemático usado
    \item Não usar símbolos desnecessários ou cujo significado não compreenda
\end{itemize}

A resolução de um exercício, em qualquer disciplina que envolva Matemática, deve ser pensada exatamente da mesma forma
que uma dissertação, a única diferença é o uso de simbologia Matemática. Ela deve conter uma introdução, 
uma sequência lógica de argumentos que levem o leitor a chegar a conclusão desejada, que é a resposta correta
do problema. Cada símbolo Matemático faz o papel de um elemento de uma frase, assim como toda fase deve 
ser completa e transmitir alguma informação coerente, toda afirmação Matemática deve ser completa e coerente.

Ao escrever uma resolução, siga a fluxo de texto natural para os leitores de Português (ou da língua em que a disciplina
estiver sendo ministrada), isso é, a sequência de passos deve seguir da esquerda para a direita e de cima para baixo.
Não é necessário indicar a sequência por setas ou outras formas gráficas. 

Também é necessário tomar cuidado com a caligrafia ao escrever as respostas na prova. O avaliador deve ser capaz
de entender o que está escrito. Use a mesma caligrafia usada para escrever a redação na prova do ENEM.

A resolução de exercícios por métodos diferentes dos apresentados no curso será considerada \textbf{errada}, 
mesmo que o método usado pelo aluno seja correto. Recomento consultar o professor sobre métodos
alternativos durante as aulas (jamais durante as provas). A razão para essa atitude é que 
ao usar um método diferente do apresentado pelo professor, pode indicar que o aluno simplesmente 
não entendeu o método apresentado.

Normalmente não é permitido o uso de calculadoras ou celulares nas aulas e provas. Os celulares por motivos óbvios
e a calculadora simplesmente não será necessária. Cálculos complexos com expressões numéricas não devem aparecer 
nas provas e caso apareçam, basta deixar a expressão indicada sem avaliar seu valor numérico.

%------------------------------------------------------------------------------%
\section{Como Estudar Matemática}

Esses são alguns passos, propostos pela Professora Sandra, que podem auxiliar no estudo de Matemática.
\begin{description}
    \item[\textcolor{blue}{Passo 1}] O início da aprendizagem, na maioria das vezes, começa em sala de aula com o(a) Professor(a). 
        Por isso, é de extrema importância que você:
        \begin{itemize}
            \item se concentre e se envolva no assunto que está sendo explicado e construído;
            \item pergunte (sem constrangimento) se algo não ficar claro no seu entendimento;
            \item desligue o celular durante a aula, e não fique se comunicando e se distraindo com redes sociais.
        \end{itemize}
    \item[\textcolor{blue}{Passo 2}] O conteúdo desenvolvido em sala de aula precisa ser amadurecido e absorvido por completo. 
        Sendo que, para que isso ocorra é crucial:
        \begin{itemize}
            \item Antes de ir para a próxima aula estude a aula anterior, reforce o conteúdo desenvolvido, 
                e sobretudo, refaça sozinho todos os exemplos feitos em sala. Isso fará com que você reforce 
                a ideia construída/utilizada em cada um deles. Além disso, faça os exercícios propostos pelo 
                professor, caso tenha sugerido. 
                \begin{center}
                    ``Matemática também se aprende por repetição''
                \end{center}
        \end{itemize}
    \item[\textcolor{blue}{Passo 3}] Encontre um parceiro/colega para estudar com você. Procure pessoas que tenham, 
        mais ou menos, o mesmo ritmo que o seu.
    \item[\textcolor{blue}{Passo 4}] Para uma boa articulação lógica do conhecimento matemático absorvido é necessário que você:
        \begin{itemize}
            \item faça ativamente um grande número de exercícios sobre o assunto. Ver passivamente o professor 
                (ou alguém) de forma presencial ou por vídeo resolvendo não é suficiente. 
            \item Lembre-se, o processo de tentativa e erro faz parte natural da resolução dos exercícios.
        \end{itemize}
    \item[\textcolor{blue}{Passo 5}] Dúvidas acumuladas geram um problema grande para sua aprendizagem, por isso:
        \begin{itemize}
            \item Ir à monitoria com frequência é extremamente importante. 
        \end{itemize}
\end{description}
Ao fazer uma avaliação sempre opte por executar inicialmente as questões que você considera mais fáceis.

%------------------------------------------------------------------------------%
\section{Recursos Online}

Existem muitos recursos \textit{online} que servem como apoio ao estudo de Matemática e Cálculo. 
Utilizar esses recursos é altamente recomendável, porém lembre-se que apenas assistir vídeos passivamente
não é suficiente para aprender Matemática, da mesma forma que, assistir atletas olímpicos não nos torna 
atletas também. 

Alguns recursos \textit{online} que podem ser uteis:
\begin{itemize}
    \item \href{http://www.wolframalpha.com}{www.wolframalpha.com} \\
         WolframAlpha é um mecanismo de conhecimento computacional que é capaz de fazer muitos cálculos.
         
    \item \href{http://pt.khanacademy.org}{pt.khanacademy.org} \\
        A Khan Academy é uma ONG educacional criada e sustentada por Sal Khan. 
        Com a missão de fornecer educação de alta qualidade para qualquer um, em qualquer lugar, 
        oferece uma coleção grátis de vídeos de matemática, medicina e saúde, economia e finanças, 
        física, química, biologia, ciência da computação, entre outras matérias.
        
    \item \href{http://www.mathway.com/pt}{www.mathway.com/pt} \\
        A Mathway fornece aos alunos ferramentas para compreender e resolver problemas matemáticos. 
        As resoluções são apresentadas passo a passo.
        
    \item \href{http://onlinemschool.com/math/assistance/equation/gaus}{onlinemschool.com/math/assistance/equation/gaus} \\
        Calculadora online que utiliza o método de Gauss para resolver um sistema linear mostrando os cálculos passo a passo.
\end{itemize}



%------------------------------------------------------------------------------%